\chapter{Tools, MCP and Skills}
\label{ch:tools}

\begin{quote}
\itshape
Giving an agent the ability to act, and the ability to be stopped before it does.
\end{quote}

\noindent
An agent without tools is a chat endpoint with a system prompt. Tools are what
make it an agent, and they are also where most of the risk lives --- a tool call
is the point at which a probabilistic system reaches into a deterministic one and
changes something.

This chapter covers four ways to give an agent capabilities: C\# methods, Model
Context Protocol servers, other agents, and skills. It also covers the mechanism
for putting a human in front of any of them.

\begin{verifybox}
As in Chapter~\ref{ch:first-agent}, the listings here are transcribed from
Microsoft documentation and published samples rather than compiled by the author
against \mafcore{}. Compile before you copy.
\end{verifybox}

% ============================================================
\section{A C\# method becomes a tool}
\label{sec:function-tools}
\index{AIFunctionFactory@\texttt{AIFunctionFactory}}

\api{AIFunctionFactory.Create} turns a method into something an agent can call.
The method can be static or instance, sync or async.

\begin{csharp}[caption={A function tool}, label={lst:tool-basic}]
using System.ComponentModel;
using Microsoft.Agents.AI;
using Microsoft.Extensions.AI;

[Description("Get the current weather for a given location.")]
static string GetWeather(
    [Description("The city name, e.g. 'Alicante' or 'Zurich'.")] string location)
    => $"The weather in {location} is cloudy with a high of 15C.";

AIAgent agent = chatClient.AsAIAgent(
    instructions: "You are a helpful assistant with access to weather information.",
    name: "WeatherAgent",
    tools: [AIFunctionFactory.Create(GetWeather)]);

Console.WriteLine(await agent.RunAsync("What is the weather in Amsterdam?"));
\end{csharp}

Note what the framework derives automatically: the tool name from the method
name, the parameter schema from the signature, and the human-readable
descriptions from the \api{[Description]} attributes.

\subsection{What the model actually sees}

This is the mental model that makes the rest of the chapter make sense. The model
does not see your method. It sees a name, a description, and a JSON schema for
the parameters --- and it chooses among them on that basis alone. It has no
access to the implementation, no type system, and no way to ask a clarifying
question about what \code{location} means.

Two consequences follow immediately.

\textbf{The description is not documentation, it is the interface.} A method named
\code{GetData} with the description ``gets data'' is genuinely ambiguous to the
model in a way it is not to you, because you have the surrounding code and it does
not.

\textbf{The schema is the contract.} If a parameter can only meaningfully take
three values, an enum communicates that and a \code{string} does not.

\subsection{The tool-calling loop}

When you call \api{RunAsync}, the framework runs a loop on your behalf: it detects
function calls in the model's response, optionally requests approval
(\S\ref{sec:approvals}), executes the functions, injects the results back into the
conversation, and repeats until the model produces a final answer or hits an
iteration limit.

\begin{warning}
\index{rate limits}
Every iteration of that loop is a separate model call. A task requiring eight tool
calls costs eight round trips, and a multi-agent orchestration multiplies it
again --- three agents at ten iterations each is thirty calls for one user
request. Under load this hits provider rate limits well before it hits anything
else. Chapter~\ref{ch:harness} discusses collapsing the loop; for now, be aware
that tool count and latency are directly proportional.
\end{warning}

% ============================================================
\section{Parameter shapes}
\label{sec:param-shapes}

Not all schemas are equally easy for a model to fill in correctly. This is the
single largest practical difference between a frontier model and a small local
one, and it is worth establishing where your chosen model stops coping.

Four shapes, in increasing order of difficulty:

\begin{csharp}[caption={Four parameter shapes}, label={lst:tool-shapes}]
// 1. Primitives -- reliable almost everywhere
[Description("Convert a temperature from Celsius to Fahrenheit.")]
static double ToFahrenheit([Description("Degrees Celsius.")] double celsius)
    => celsius * 9 / 5 + 32;

// 2. Constrained values -- an enum tells the model what is allowed
enum Unit { Metric, Imperial }

[Description("Get the weather in the requested unit system.")]
static string GetWeather(string location, Unit unit) => /* ... */;

// 3. A structured argument
record Booking(
    [property: Description("ISO-8601 date, e.g. 2026-08-14.")] string Date,
    [property: Description("Number of guests, 1 to 8.")] int Guests);

[Description("Reserve a table.")]
static string Reserve(Booking booking) => /* ... */;

// 4. Collections of structured arguments -- the hardest case
[Description("Reserve several tables in one request.")]
static string ReserveMany(IReadOnlyList<Booking> bookings) => /* ... */;
\end{csharp}

\begin{note}
Prefer several narrow tools over one tool with a discriminated payload. It is
tempting to write \code{ExecuteAction(string actionType, string payloadJson)} and
let the model fill in the details. Models select poorly against that shape,
because the description cannot say what \code{payloadJson} should contain for each
\code{actionType}. Three separate tools with three specific schemas outperform it
consistently.
\end{note}

\subsection{Async tools and failure}

Async methods work directly --- return \code{Task<T>} and register the same way.
Failure needs more thought.

An exception thrown inside a tool has to become something the model can read,
because the model is the caller. Returning a structured failure that says what
went wrong is usually better than throwing, since the model can then retry with
different arguments or tell the user what it could not do. Throwing is
appropriate when the failure is not the model's to recover from --- a
misconfiguration, an authentication failure, a bug.

The distinction is worth being deliberate about: \emph{recoverable by choosing
different arguments} means return a message; \emph{not recoverable by the model}
means throw.

% ============================================================
\section{Measuring selection accuracy}
\label{sec:tool-accuracy}

The claim that descriptions matter is repeated everywhere and quantified almost
nowhere. Quantify it.

\begin{exercisebox}
Take one tool and write three variants of its description:

\begin{enumerate}[leftmargin=1.4em]
\item \textbf{Minimal} --- ``Gets weather.''
\item \textbf{Specific} --- what it returns, in what units, for what input.
\item \textbf{Specific with negative guidance} --- as above, plus what it is
      \emph{not} for, and when the model should not call it.
\end{enumerate}

Register it alongside two plausible distractor tools. Run twenty identical
prompts against each variant and record how often the correct tool was chosen
with correct arguments.

Then repeat the whole thing for each parameter shape in \S\ref{sec:param-shapes},
against both a local and a hosted model. Record the results in
Appendix~\ref{app:providers}.
\end{exercisebox}

This is the most transferable measurement in the book. The absolute numbers are
model-specific and will age; the \emph{shape} of the result --- how sharply
accuracy falls off with schema complexity on a small model --- is what tells you
whether a cheap provider is viable for a given workload.

\begin{warning}
Before concluding that a description is inadequate, rule out the two mechanical
causes: the model may not support function calling at all, and some chat clients
require the function-invocation middleware to be added explicitly. See
Appendix~\ref{app:troubleshooting}, \S\ref{sec:tb-no-tool-call}, and work through
it in order. Rewriting prompts to fix a missing \api{UseFunctionInvocation} call
is a memorable way to lose an afternoon.
\end{warning}

% ============================================================
\section{Human-in-the-loop approval}
\label{sec:approvals}
\index{approval!function}
\index{human-in-the-loop}

By the time an agent has tools, the question of what happens when it decides to
take an irreversible action becomes practical rather than theoretical. Cancelling
an order, charging a card, deleting a record --- these are not decisions to leave
to a model, however confident the prose around the call sounds.

The framework's answer is to wrap the tool so that it cannot execute without an
explicit decision.

\subsection{Requiring approval}

Sensitive tools are wrapped with \api{ApprovalRequiredAIFunction}, which comes
from \pkg{Microsoft.Extensions.AI} rather than from Agent Framework itself.

\subsection{The request}

An agent run against an approval-required tool does not complete with an answer.
It completes with a request for input:

\begin{csharp}[caption={Detecting an approval request}, label={lst:approval-detect}]
AgentSession session = await agent.CreateSessionAsync();
AgentResponse response = await agent.RunAsync("What is the weather in Amsterdam?", session);

var approvalRequests = response.Messages
    .SelectMany(x => x.Contents)
    .OfType<FunctionApprovalRequestContent>()
    .ToList();

foreach (var request in approvalRequests)
{
    Console.WriteLine($"Approval required for '{request.FunctionCall.Name}'");
    // request.FunctionCall.Arguments carries what the model wants to pass.
}
\end{csharp}

\api{FunctionApprovalRequestContent.FunctionCall} carries the name and arguments,
which is exactly what you show a human: \emph{the agent wants to call
\code{CancelOrder} with \code{orderId = 4471}. Allow?}

\subsection{The response}

\begin{csharp}[caption={Approving, and continuing the run}, label={lst:approval-respond}]
FunctionApprovalRequestContent requestContent = approvalRequests.First();

// The user has decided. Create the matching response.
var approval = requestContent.CreateResponse(approved: true);

// Pass it back as part of a new run on the same session.
response = await agent.RunAsync(new ChatMessage(ChatRole.User, [approval]), session);
\end{csharp}

\begin{warning}
\textbf{Keep checking until there are none left.} One run can produce several
approval requests, and answering one can produce more on the next turn. The rule
is: after every agent run, check for \api{FunctionApprovalRequestContent} and
loop until none remain. Code that checks once handles the demo and fails the
third time a user asks for something complicated.
\end{warning}

\needscreenshot{vs-debugger-approval-request}%
{An approval request, mid-run}%
{\vsver{} stopped at a breakpoint on the \texttt{foreach} in
Listing~\ref{lst:approval-detect}. Capture the \emph{Locals} or \emph{Watch}
window with \texttt{request} expanded to show \texttt{FunctionCall.Name} and
\texttt{FunctionCall.Arguments} with real values visible --- the point of the
figure is that the reader sees exactly what information is available to show a
human. Widen the pane so the arguments are not truncated.}

\subsection{Approvals inside workflows}

The same mechanism carries into Part~II. When an agent inside an orchestration
calls a tool wrapped with \api{ApprovalRequiredAIFunction}, the workflow pauses
and emits a \api{RequestInfoEvent} carrying a \api{ToolApprovalRequestContent}.
An external system inspects it, approves or rejects, and the workflow resumes:

\begin{csharp}[caption={The same approval, from inside a workflow}, label={lst:approval-workflow}]
await run.SendResponseAsync(
    e.Request.CreateResponse(approvalRequest.CreateResponse(approved: true)));
\end{csharp}

Sequential orchestration supports this with no additional configuration.
Chapter~\ref{ch:workflows-advanced} covers the request-and-response model
properly; the point here is that approval is one mechanism, not two.

\begin{note}
Design approvals around \emph{irreversibility}, not around sensitivity. A tool
that reads a customer record is sensitive but recoverable; a tool that emails the
customer is not. The second needs a gate and the first probably does not, and
gating everything trains operators to click through without reading --- which is
worse than no gate at all.
\end{note}

% ============================================================
\section{Model Context Protocol}
\label{sec:mcp}
\index{MCP}

MCP is an open standard letting agents call tools hosted by external servers,
regardless of what language those servers are written in. Agent Framework
integrates with it through the official .NET SDK rather than a bespoke
abstraction.

\begin{center}
\small
\begin{tabularx}{\linewidth}{@{}lX@{}}
\toprule
\textbf{Package} & \textbf{Use} \\
\midrule
\pkg{ModelContextProtocol} & Most client and stdio server scenarios. \\
\pkg{ModelContextProtocol.Core} & Low-level or client-only, fewer dependencies. \\
\pkg{ModelContextProtocol.AspNetCore} & Servers hosted over HTTP in ASP.NET Core. \\
\bottomrule
\end{tabularx}
\end{center}

\begin{versionbox}
These packages reached stable 1.4.0 in \mafdate{}, on their own release cadence
entirely separate from Agent Framework's. Older material adds
\code{-{}-prerelease}; that is no longer required. The MCP specification is still
evolving, so track the SDK repository rather than assuming the API is settled.
\end{versionbox}

\subsection{Consuming a server}
\label{sec:mcp-client}

Three steps: open a connection, discover what the server offers, hand the
discovered tools to the agent.

\begin{csharp}[caption={An agent backed by an external MCP server}, label={lst:mcp-client}]
using Microsoft.Agents.AI;
using Microsoft.Extensions.AI;
using ModelContextProtocol.Client;

// 1. Connect. StdioClientTransport launches a local server process.
await using var mcpClient = await McpClient.CreateAsync(new StdioClientTransport(new()
{
    Name = "GitHub",
    Command = "npx",
    Arguments = ["-y", "@modelcontextprotocol/server-github"]
}));

// 2. Discover. Nothing is hardcoded.
IList<McpClientTool> mcpTools = await mcpClient.ListToolsAsync();

// 3. Register. McpClientTool implements AITool, so it goes straight in.
AIAgent agent = chatClient.AsAIAgent(
    instructions: "You answer GitHub repository questions using the provided tools.",
    name: "GitHubAgent",
    tools: [.. mcpTools]);
\end{csharp}

\begin{note}
\textbf{\api{McpClient.CreateAsync}, not \api{McpClientFactory.CreateAsync}.}
Both appear in published examples. The factory type is the older spelling; every
current sample in the Agent Framework repository uses \api{McpClient.CreateAsync}
with the transport options supplied as an object initialiser. The three-step
shape --- connect, discover, register --- is stable across both.
\end{note}

Step two is the reason MCP is interesting. You do not write bindings for the
server's tools; you ask it what it has. A server that gains a tool next month
requires no change on your side.

\begin{warning}
\index{MCP!disposal}
\textbf{Use \code{await using}, not \code{using}.} An MCP client holds a child
process or a network socket open, and disposal is asynchronous. Getting this
wrong leaks processes --- invisibly in a console sample, expensively in a
long-running service that creates a client per request.
\end{warning}

\subsection{Transports}

\api{StdioClientTransport} launches a local process and speaks over its standard
input and output. It suits local, desktop and IDE-hosted tooling.
\api{SseClientTransport} connects to a remote server over HTTP. The choice is
mostly about where the server runs and who owns it.

\needscreenshot{vs-debugger-mcp-tools}%
{Tools discovered from an MCP server at run time}%
{\vsver{} stopped at a breakpoint immediately after \texttt{ListToolsAsync()} in
Listing~\ref{lst:mcp-client}. Capture the \emph{Locals} window with the
\texttt{mcpTools} collection expanded so that several discovered tool names and
their descriptions are readable. The figure exists to show that discovery is
dynamic --- so the count and the names being visible matters more than any
individual entry.}

\subsection{Exposing an agent as an MCP tool}
\label{sec:mcp-server}

The relationship runs both ways. An agent can be published as an MCP tool that
anything speaking the protocol can call --- including agents written in other
languages, and including tooling that is not an agent at all.

\begin{csharp}[caption={Publishing an agent over MCP}, label={lst:mcp-server}]
using Microsoft.Agents.AI;
using Microsoft.Extensions.DependencyInjection;
using Microsoft.Extensions.Hosting;
using ModelContextProtocol.Server;

AIAgent nutritionAgent = chatClient.AsAIAgent(
    instructions: "You can fetch recipes and nutrition data.",
    name: "NutritionAgent",
    description: "Searches for high-protein recipes and nutrition information.",
    tools: [AIFunctionFactory.Create(SearchRecipesAsync)]);

// The agent's name and description become the MCP tool's name and description.
McpServerTool tool = McpServerTool.Create(nutritionAgent.AsAIFunction());

HostApplicationBuilder builder = Host.CreateEmptyApplicationBuilder(settings: null);
builder.Services
    .AddMcpServer()
    .WithStdioServerTransport()
    .WithTools([tool]);

await builder.Build().RunAsync();
\end{csharp}

Note the \code{description} argument on the agent. In
Chapter~\ref{ch:first-agent} the \code{name} was described as functional rather
than decorative; here both name and description become the interface other
systems select against. This is the second place where those strings stop being
cosmetic, and it will not be the last.

\subsection{Security}

MCP moves the trust boundary. Three things follow.

\textbf{Server credentials belong in the environment, not in your source.} The
GitHub server in Listing~\ref{lst:mcp-client} reads a personal access token from
the process environment. Set it in the shell or through user secrets; never
inline it.

\textbf{Tool results are untrusted input.} An MCP server returns text that goes
straight into the model's context. If the server is not yours --- and often it is
not --- that text can contain instructions. This is prompt injection with a
supply chain attached to it, and the defences are the ones you would apply to any
untrusted input: constrain what the agent can do with the result, and gate the
consequential actions behind approval (\S\ref{sec:approvals}).

\textbf{Discovery is dynamic, which cuts both ways.} A server that gains a tool
next month gains it in your agent too, without a code review. For third-party
servers, pin the version.

% ============================================================
\section{Agents as tools}
\label{sec:agents-as-tools}

\api{AsAIFunction} works outside MCP as well. Any agent can become a tool for
another agent:

\begin{csharp}[caption={An orchestrator delegating to specialists}, label={lst:agent-as-tool}]
AIAgent orchestrator = chatClient.AsAIAgent(
    instructions: """
        You are an assistant with access to specialist experts.
        Delegate to the most appropriate expert using your tools,
        then combine their answers into one coherent response.
        """,
    name: "Orchestrator",
    tools: [historyAgent.AsAIFunction(), scienceAgent.AsAIFunction()]);
\end{csharp}

This is the simplest multi-agent pattern in the framework, and it is worth
understanding what it is not. The specialists here are \emph{tools}: they are
called, they return, control comes back. They do not take over the conversation,
they cannot hand off to each other, and the orchestrator remains responsible for
the final answer.

Chapter~\ref{ch:orchestration} covers handoff, where control genuinely transfers.
The distinction is delegation versus transfer, and choosing wrongly is a common
source of multi-agent systems that are harder to reason about than they need to
be. If the caller should own the final answer, use a tool.

% ============================================================
\section{Agent Skills}
\label{sec:skills}
\index{skills}

Skills package instructions, resources and scripts on a filesystem so an agent
can discover and load them on demand, rather than carrying everything in its
system prompt.

The skills provider exposes exactly three tools to the agent:

\begin{center}
\small
\begin{tabularx}{\linewidth}{@{}lX@{}}
\toprule
\textbf{Tool} & \textbf{Does} \\
\midrule
\code{load\_skill} & Loads a skill's instructions into context. \\
\code{read\_skill\_resource} & Fetches a resource bundled with a skill. \\
\code{run\_skill\_script} & Executes a script bundled with a skill. \\
\bottomrule
\end{tabularx}
\end{center}

\textbf{All three require approval by default.} Nothing loads and nothing executes
without oversight, and you relax that selectively for operations you trust. Given
that the third tool executes code, this default is the correct one and you should
think hard before changing it.

\subsection{Where scripts run}

The execution model has a split that matters:

\begin{itemize}[leftmargin=1.4em]
\item \textbf{Class-based and code-defined skill scripts run in-process.} They are
      your code; the framework calls them.
\item \textbf{File-based scripts are delegated to a runner you provide.} The
      framework does not execute them for you.
\end{itemize}

That second point is a design decision, not an omission. Sandboxing, resource
limits and audit logging for arbitrary file-based scripts are yours to implement,
because only you know what isolation your environment requires.

\begin{warning}
If you implement a file-based script runner, treat it as a security boundary from
the first line. It executes code selected by a model on the basis of text that may
have come from an untrusted tool result. The combination of dynamic skill
discovery and script execution is the most consequential capability in this
chapter --- and, unlike the approval gate, nothing about it is safe by default
once you have written the runner.
\end{warning}

\begin{versionbox}
Agent Skills for .NET is a recent release. Check the current documentation for
the provider's configuration surface, including the filtering options for
controlling which skills an agent can discover.
\end{versionbox}

% ============================================================
\section{Summary}

A C\# method becomes a tool through \api{AIFunctionFactory.Create}, with the name
taken from the method, the schema from the signature, and the descriptions from
\api{[Description]}. The model sees only those three things, which is why
descriptions are an interface rather than documentation, and why several narrow
tools beat one tool with a discriminated payload.

Selection accuracy degrades with schema complexity, sharply so on smaller models.
Measure it rather than assuming it --- but rule out the mechanical causes first.

Sensitive tools are wrapped with \api{ApprovalRequiredAIFunction}. The run
returns \api{FunctionApprovalRequestContent} instead of an answer; you show the
function name and arguments to a human, call \api{CreateResponse}, and continue
--- looping until no requests remain. The same mechanism appears inside workflows
as \api{RequestInfoEvent}. Gate on irreversibility, not on sensitivity.

MCP is connect, discover, register --- and \code{await using}, not \code{using}.
Discovery being dynamic is the feature and also the reason to pin third-party
servers. Tool results are untrusted input.

\api{AsAIFunction} turns an agent into a tool, for MCP or for another agent. That
is delegation, and it is a different thing from the handoff in
Chapter~\ref{ch:orchestration}.

Skills add filesystem-discovered instructions, resources and scripts, gated by
approval on all three operations by default. File-based script execution is
delegated to a runner you write, which makes the isolation your responsibility.

\begin{exercisebox}
Before Chapter~\ref{ch:memory}:

\begin{enumerate}[leftmargin=1.4em]
\item Register three tools of increasing schema complexity and confirm your local
      model calls each one correctly. Note where it stops coping.
\item Run the description experiment in \S\ref{sec:tool-accuracy} and record the
      numbers in Appendix~\ref{app:providers}.
\item Wrap a destructive tool with \api{ApprovalRequiredAIFunction} and implement
      the full approval loop --- including the case where one run produces two
      requests.
\item Connect an agent to an existing MCP server, then write a small MCP server of
      your own and consume it. Exposing something you already have --- a
      diagnostic query, a build status, a log search --- is more instructive than
      a toy.
\item Take one tool from step~1 and expose it through your MCP server instead.
      Note what changed in the agent code. It should be almost nothing.
\end{enumerate}
\end{exercisebox}
